\conclusion
Au terme de ce projet, réalisé au sein d’Inetum Offshore Maroc, l’objectif majeur de modernisation du Système d’Information de l'opérateur telecom a été pleinement atteint pour le périmètre du Lot 1. Ce stage a représenté une immersion complète dans un projet de transformation numérique d'envergure industrielle, où l'enjeu était de centraliser l'intégralité du patrimoine réseau et client au sein d'une « Source Unique de Vérité » via la solution Ni2 Unified Inventory.\\

Le travail que j'ai mené s'est articulé autour de la convergence entre les infrastructures techniques et les référentiels commerciaux. Sur le volet technique, ma contribution a permis de stabiliser l'inventaire des réseaux cuivre et fibre issu du système GAIA. Concernant le volet commercial, j'ai piloté une démarche de rétro-ingénierie méthodique sur la base de données CRM. Ce travail d'investigation a été indispensable pour décoder les règles de gestion métier et assurer une réconciliation fluide des données clients au sein du nouvel outil.\\

Toutefois, le projet de l'operateur teleco est une aventure technologique qui se poursuit au-delà de la rédaction de ce document. Les travaux réalisés jusqu'ici posent les fondations nécessaires au déploiement du Jalon 5, lequel constitue mon prochain objectif opérationnel. Cette phase à venir sera notamment consacrée à la finalisation des tests d'intégration de bout en bout entre le CRM, la Façade et Ni2, afin de valider la fluidité des processus métiers dans leur globalité.\\

Sur le plan personnel, ce stage est un véritable catalyseur pour mon profil d'ingénieure en Réseaux et Systèmes Informatiques. Il m'a permis de développer une solide expertise en SQL avancé et en gestion de systèmes d'information massifs, tout en m'apprenant à évoluer avec agilité dans un environnement de production international.\\

En conclusion, si ce mémoire documente une étape clé de la migration, il témoigne surtout d'une dynamique de travail continue. La fiabilisation de l'inventaire technique et la réconciliation commerciale déjà accomplies représentent un succès intermédiaire décisif, ouvrant la voie à la livraison finale du Lot 1 et, à terme, à la modernisation complète des services de l'opérateur.
